\chapter{Human Chorionic Gonadotropin (hCG) Test}

\section*{Definition and Overview}
The human chorionic gonadotropin (hCG) test measures the level of hCG, a hormone produced by the placenta shortly after a fertilised egg implants in the uterus. It is best known as the pregnancy test, and it is also used to monitor pregnancy health, diagnose and follow ectopic pregnancy, and manage certain cancers that produce hCG. The test is available as a urine test, which detects pregnancy, and as a blood test, which measures the level precisely. hCG levels rise rapidly in early pregnancy and follow a characteristic pattern.

\section*{Epidemiology}
Pregnancy testing is among the most common tests performed in the world. Around 300,000 babies are born in Australia each year, and countless pregnancy tests are used by women at home and in health settings. hCG measurement is also essential in the management of ectopic pregnancy, which affects around 2\% of pregnancies and is a leading cause of early pregnancy death, and in the follow-up of gestational trophoblastic disease and some testicular cancers.

\section*{Aetiology and Pathophysiology}
\begin{itemize}
    \item \textbf{Production:} hCG is produced by the trophoblast cells of the developing placenta
    \item \textbf{Detection:} urine tests detect hCG about 12--14 days after conception, around the time of the missed period
    \item \textbf{Rising levels:} hCG roughly doubles every 48--72 hours in early normal pregnancy
    \item \textbf{Peak and decline:} levels peak around 8--10 weeks and then fall to a lower plateau
    \item \textbf{Ectopic pregnancy:} hCG rises more slowly than in a normal pregnancy
    \item \textbf{Other causes:} some tumours, including gestational trophoblastic disease and testicular germ cell tumours, produce hCG
    \item \textbf{Pathophysiology:} hCG maintains the corpus luteum, which produces progesterone to support early pregnancy
\end{itemize}

\section*{Clinical Features}
\begin{itemize}
    \item A positive pregnancy test is usually the first indication of pregnancy
    \item Home urine tests are highly accurate when used correctly
    \item Blood tests measure exact levels and are used in early pregnancy complications
    \item Falling or slowly rising hCG may indicate miscarriage or ectopic pregnancy
    \item Symptoms such as pain or bleeding with a positive test require urgent assessment
    \item hCG is also measured after treatment of trophoblastic disease and some cancers
\end{itemize}

\section*{Differential Diagnosis}
\begin{itemize}
    \item A positive test can occur in normal pregnancy, ectopic pregnancy, miscarriage, and rarely trophoblastic disease or hCG-producing tumours
    \item A false positive is uncommon but possible with some medical conditions or test misuse
    \item A false negative occurs with very early testing or dilute urine
    \item Quantitative testing and ultrasound clarify the situation
    \item The clinical picture, not the hCG level alone, guides management
\end{itemize}

\section*{When to Seek Medical Care}
A positive home pregnancy test should be confirmed with a health professional, and antenatal care should begin early. Any pregnancy with abdominal pain, bleeding, or dizziness requires urgent assessment to exclude ectopic pregnancy.

\textbf{Call 000 (or local emergency services) for:}
\begin{itemize}
    \item Severe abdominal or pelvic pain with a positive pregnancy test
    \item Heavy vaginal bleeding in early pregnancy
    \item Shoulder tip pain or dizziness with bleeding
    \item Collapse or fainting
    \item Any symptoms suggesting ectopic pregnancy rupture
\end{itemize}

\section*{Diagnosis and Investigations}
\begin{itemize}
    \item \textbf{Urine test:} rapid detection of hCG for pregnancy diagnosis
    \item \textbf{Quantitative blood test:} measures the exact hCG level
    \item \textbf{Serial testing:} repeated levels 48 hours apart assess pregnancy viability
    \item \textbf{Ultrasound:} confirms the location of the pregnancy
    \item \textbf{Other tumour markers:} where cancer is suspected
    \item \textbf{Follow-up testing:} hCG is monitored to undetectable levels after treatment of trophoblastic disease
\end{itemize}

\section*{Management}
\begin{itemize}
    \item \textbf{Normal pregnancy:} confirmation and early antenatal care
    \item \textbf{Ectopic pregnancy:} urgent treatment with medication or surgery
    \item \textbf{Miscarriage:} support and management of bleeding
    \item \textbf{Trophoblastic disease:} surgical evacuation and hCG monitoring
    \item \textbf{Cancer management:} hCG as a tumour marker in treatment and follow-up
    \item \textbf{Explanation:} results are explained clearly with appropriate support
\end{itemize}

\section*{Complications}
\begin{itemize}
    \item Missed ectopic pregnancy with rupture and life-threatening bleeding
    \item Undiagnosed miscarriage causing distress and bleeding
    \item False reassurance from negative tests taken too early
    \item Anxiety from uncertain or discordant results
    \item Delayed cancer diagnosis in rare hCG-producing tumours
\end{itemize}

\section*{Prognosis and Outcomes}
The hCG test is highly reliable and has transformed early pregnancy care. Most positive tests lead to healthy pregnancies, and ectopic pregnancies are detected and treated before rupture when women present promptly. Trophoblastic disease is curable with monitoring of hCG levels. The combination of hCG testing and ultrasound allows accurate, early diagnosis of pregnancy complications, improving outcomes for both mother and baby.

\section*{Prevention and Self-Care}
\begin{itemize}
    \item Test at the right time, ideally after the missed period, using morning urine
    \item Confirm a positive home test with a doctor
    \item Begin antenatal care early in pregnancy
    \item Seek urgent review for pain or bleeding in early pregnancy
    \item Attend follow-up testing as advised
    \item Use contraception as planned when pregnancy is not intended
    \item Understand that one negative test does not rule out pregnancy if testing too early
\end{itemize}

\section*{Sources and Further Reading}
The following reputable sources provide further information for healthcare students and patients seeking reliable, current advice about hCG testing.

\begin{itemize}
    \item Healthdirect. \textit{Pregnancy testing (hCG test).} https://www.healthdirect.gov.au/pregnancy-test
    \item Royal Women's Hospital. \textit{Early pregnancy complications and hCG.} https://www.thewomens.org.au/
    \item Better Health Channel. \textit{Pregnancy tests.} https://www.betterhealth.vic.gov.au/
    \item NHS. \textit{Pregnancy test: when to take one.} https://www.nhs.uk/pregnancy/
    \item Mayo Clinic. \textit{Pregnancy tests and hCG.} https://www.mayoclinic.org/
    \item MedlinePlus. \textit{hCG blood test.} https://medlineplus.gov/ency/article/003509.htm
\end{itemize}
